\subsection*{Trigger C++ Template Instantiation}

It is possible for Croqtile code to make function call to a template function. For example:

\begin{lstlisting}[language=croqtile,caption={A templated device helper}]
template<int M, int N, int K>
__device__ void matmul_kernel(int *lhs, int* rhs, int* output) {}
__co__ auto matmul(s32 [96, 72] lhs, s32 [72, 24] rhs) {
  s32 [lhs.span(0), rhs.span(1)] output;
  parallel {p, q} by [6, 6] : block {
    shared s32 [output.span / {#p, #q}] buffer;   // shape: [16, 4]
    lhs_load = dma.copy lhs.chunkat(p, _) => shared;
    rhs_load = dma.copy rhs.chunkat(_, q) => shared;
    call matmul_kernel<buffer.span(0), buffer.span(1), 72>(lhs_load.data, rhs_load.data, buffer);
    dma.copy buffer => output.chunkat(p, q);
  }
  return output;
}
\end{lstlisting}

The call supplies three constant template arguments to \co@matmul_kernel@. The target compiler instantiates the corresponding specialization, just as it would for an ordinary C++ template call.

Template arguments must be evaluable as constants during Croqtile compilation. An input-dependent value cannot select a C++ specialization at runtime:

\begin{lstlisting}[language=croqtile,caption={A call that instantiates it}]
__co__ void foo(int M) {
  parallel p by 1 {
    call bar<M>();  // warning: 'M' cannot be evaluated at compile time
  }
}
\end{lstlisting}

The Croqtile compiler attempts to evaluate template arguments at compile time. If it cannot determine a required constant, the compiler reports a warning that the argument cannot be evaluated at compile time.
